\section{Tensors}

This section will be rather short, assuming a basic familiarity with tensor products.
The following definition is a brief overview of notation.

\bdefn

    We denote the tensor product of modules by $A\otimes B$.
    We denote the dual of an $R$-module $A$ by $A^*=\hom_R(A,R)$.

\edefn

\bnote

    Primitive elements of the tensor product $A\otimes B$ are written $a\otimes b$; all other elements are sums of these.

    The tensor product is associative (in a natural way), and we denote the tensor of $n$ modules without parentheses.
    In an $n$-fold tensor product, we write the primitive elements without parentheses.

    If $B_i=(b^i_1,\dots,b^i_{n_i})$ is a basis for each $V_1$, then
    $$ B = \set{b^1_{j_1}\otimes b^n_{j_n}}[b^i_{j_i}\in B_i] \subseteq V_1\otimes\cdots\otimes V_n $$
    is a basis for the tensor product.
    In particular the dimension of a tensor product of vector spaces is the product of the dimensions of the vector spaces.

    If $A$ is an $R$-module, then $A\otimes R\cong A$ naturally.
    This natural isomorphism is given by mapping $a\otimes r\mapsto ra$.

    The tensor product is functorial: if $F_i\colon A_i\to B_i$ are maps for $i=1,2$ then define
    $F_1\otimes F_2\colon A_1\otimes A_2\to B_1\otimes B_2$ by
    $$ (F_1\otimes F_2)(a_1\otimes a_2) = (F_1(a_1))\otimes(F_2(a_2)) $$
    and extending by linearity.
    This is well-defined.

    In the case of linear functionals $\phi_i\colon A_i\to\bR$ we have that $\bR\otimes\bR\cong\bR$, and this map is naturally isomorphic to
    $$ (\phi_1\otimes\phi_2)(a_1\otimes a_2) = \phi_1(a_1)\phi_2(a_2) . $$

    The universal property of the tensor product is that
    $$ \hom(V_1\otimes\cdots\otimes V_n,A) \cong L(V_1,\dots,V_n;A) $$
    naturally, where the right side is the space of multilinear maps $V_1\otimes\cdots\otimes V_n\to A$.

\enote

\bnote

    For a field $k$ and a finite-dimensional $k$-linear space $V$, we know that $\dim V^*=\dim V$.
    For a basis $B=(v_1,\dots,v_n)$ we define its {\emphcolor dual basis} $B^*=(v_1^*,\dots,v_n^*)$ in $V^*$ as the linear functionals
    defined by
    $$ v_i^*(v_j) = \delta_{ij} , $$
    and extending by linearity.
    These are linearly independent, and thus must be a basis considering dimensions.

    There is a natural map from a space to its double dual, $V\to V^{**}$.
    For $v\in V$ define $\hat v\in V^{**}$ by $\hat v(\phi)=\phi(v)$ for all $\phi\in V^*$.
    This is injective, and considering dimensions is an isomorphism.
    It is natural, and thus for finite-dimensional spaces $V$ and $V^{**}$ are naturally isomorphic.

    For infinite-dimensional spaces, $\dim V^*$ is strictly greater than $\dim V$, so in particular $V$ is naturally a proper subspace of
    $V^{**}$, but they are not isomorphic.

\enote

\bprop

    Let $V_1,\dots,V_n$ be vector spaces, then there is a natural isomorphism
    $$ (V_1\otimes\cdots\otimes V_n)^* \cong V_1^*\otimes\cdots\otimes V_n^* . $$

\eprop

\bproof

    They have the same dimension, and thus we need only find an injective map.
    We define a map $V_1^*\otimes\cdots\otimes V_n^*$ to the dual space on the left by mapping $\phi_1\otimes\cdots\otimes\phi_n$ to the
    literal tensor product of maps:
    $$ (\phi_1\otimes\cdots\otimes\phi_n)(a_1,\dots,a_n) = \phi_1(a_1)\cdots\phi_n(a_n) . $$
    If we choose a dual basis for each $V_i^*$, we see that it maps to a linearly independent set, and thus this is injective.
    \qed

\eproof

In particular, we then have the natural isomorphisms
$$ V_1\otimes\cdots\otimes V_n \cong (V_1\otimes\cdots\otimes V_n)^{**} \cong (V_1^*\otimes\cdots\otimes V_n^*)^* . $$
Thus,
$$ V_1\otimes\cdots\otimes V_n \cong \hom(V_1^*\otimes\cdots\otimes V_n,\bR) \cong L(V_1^*,\dots,V_n^*,\bR) . $$
So we can identify the tensor product $V_1\otimes\cdots\otimes V_n$ with multilinear maps $V_1^*\times\cdots\times V_n^*\to\bR$ naturally.
Similarly, $V_1^*\otimes\cdots\otimes V_n^*$ can naturally be identified with multilinear maps $V_1\times\cdots\times V_n\to\bR$.

\bdefn

    For a space $A$, denote by $T^k(A)$ the $k$-fold tensor product of $A$ with itself:
    $$ T^k(A) = A\otimes\cdots\otimes A . $$
    The empty tensor product is taken to be the ambient ring, so $T^0(A)=R$.

\edefn

\bdefn

    Let $V$ be a finite-dimensional real vector space.
    A {\emphcolor covariant $k$-tensor} for $k>0$ is an element of $V^*\otimes\cdots\otimes V^*$.
    Equivalently, it is a multilinear map $V\times\cdots\times V\to\bR$.
    If $\alpha$ is a covariant $k$-tensor for some $k$, it is called a {\emphcolor covariant tensor} and its {\emphcolor rank} is $k$.
    The space of covariant $k$-tensors is $T^k(V^*)$.

    A {\emphcolor contravariant $k$-tensor} for $k>0$ is an element of $V\otimes\cdots\otimes V=T^k(V)$.
    Equivalently, a multilinear map $V^*\times\cdots\times V^*\to\bR$.
    The rank is defined similarly.

    For $k,\ell\neq0$ nonnegative, we define a {\emphcolor tensor of type $(k,\ell)$} to be an element of
    $$ T^{(k,\ell)}(V) = V\otimes\cdots\otimes V\otimes V^*\otimes\cdots\otimes V^* = T^k(V) \otimes T^\ell(V^*) . $$

\edefn

Note that
$$ T^{(0,k)}(V) = T^k(V^*),\qquad T^{(\ell,0)}(V) = T^\ell(V) . $$

\bdefn

    We define an action of $S_k$ on $T^k(V)$, viewed as multilinear maps, by
    $$ {}^\sigma\alpha(v_1,\dots,v_n) = \alpha(v_{\sigma(1)},\dots,v_{\sigma(n)}) $$
    for $\sigma\in S_k,\alpha\in T^k(V)$.
    This is indeed a group action, moreover it preserves the linear structure of $T^k(V)$ in that
    $$ {}^\sigma(\alpha+\beta) = {}^\sigma\alpha + {}^\sigma\beta,\qquad {}^\sigma(c\alpha) = c{}^\sigma\alpha . $$

\edefn

Viewing $T^k(V)$ as the $k$-fold tensor product of $V^*$, ${}^\sigma$ acts on primitive elements by
$$ \sigma(v_1\otimes\cdots\otimes v_k) = v_{\sigma^{-1}(1)}\otimes\cdots\otimes v_{\sigma^{-1}(k)} , $$
and extending by linearity.

\bdefn

    A covariant $k$-tensor $\alpha\in T^k(V^*)$ is {\emphcolor symmetric} if swapping its inputs does not change its value.
    Equivalently it has a trivial orbit under the above action: ${}^\sigma\alpha=\alpha$ for all $\sigma\in S_k$ (since transpositions
    generate $S_k$).
    Denote the space of symmetric covariant $k$-tensors by $\Sigma^k(V^*)\subseteq T^k(V^*)$.

    Similarly a covariant $k$-tensor $\alpha\in T^k(V^*)$ is {\emphcolor antisymmetric} if swapping its inputs swaps the sign of its value.
    Equivalently if ${}^\sigma\alpha=\sgn\sigma\alpha$ for all $\sigma\in S_k$.
    Denote the space of alternating covariant $k$-tensors by $\Lambda^k(V^*)\subseteq T^k(V^*)$.
    Antisymmetric tensors are also called {\emphcolor exterior forms} or {\emphcolor multi-covectors}.

\edefn

Define the map $\sym\colon T^k(V^*)\to\Sigma^k(V^*)$ by
$$ \sym\alpha = \frac1{k!}\sum_{\sigma\in S_k}{}^\sigma\alpha . $$
This map is called {\emphcolor symmetrization}, and it is a projection onto $\Sigma^k(V^*)$.
Indeed,
$$ {}^\tau\sym\alpha = \frac1{k!}\sum_{\sigma\in S_k}{}^{\tau\sigma}\alpha = \sym\alpha $$
so that $\sym$ is well-defined, and $\sym\alpha=\alpha$ iff $\alpha$ is symmetric.
Clearly the equality implies $\alpha$'s symmetry since $\sym\alpha$ is symmetric, and if $\alpha$ is symmetric then
$$ \sym\alpha = \frac1{k!}\sum_{\sigma\in S_k}\alpha = \alpha . $$

\bdefn

    If $\alpha\in\Sigma^k(V^*)$ and $\beta\in\Sigma^\ell(V^*)$ are symmetric, then their tensor $\alpha\otimes\beta$ is not symmetric in general.
    Define their {\emphcolor symmetric product} to be the symmetrification of their tensor product:
    $$ \alpha\beta = \sym(\alpha\otimes\beta) . $$
    Explicitly,
    $$ \alpha\beta(v_1,\dots,v_{k+\ell}) =
    \frac1{(k+\ell)!}\sum_{\sigma\in S_{k+\ell}}\alpha(v_{\sigma1},\dots,v_{\sigma k})\beta(v_{\sigma(k+1)},\dots,v_{\sigma\ell}) . $$

\edefn

\bprop

    The symmetric product satisfies the following:
    \benum
        \item it is bilinear,
        \item it is commutative,
        \item it is associative.
    \eenum

\eprop

We will not prove this.

We can also define the {\emphcolor alternation} of a tensor by
$$ \alt\alpha = \frac1{k!}\sum_{\sigma\in S_k}\sgn\sigma{}^\sigma\alpha . $$
This is indeed antisymmetric,
$$ {}^\tau\alt\alpha = \frac1{k!}\sum_\sigma\sgn\sigma{}^{\tau\sigma}\alpha = \sgn\tau\alt\alpha . $$
If $\alpha$ is already antisymmetric then
$$ \alt\alpha = \frac1{k!}\sum_\sigma\sgn\sigma{}^\sigma\alpha = \frac1{k!}\sum_\sigma(\sgn\sigma)^2\alpha = \alpha $$
so this is also a projection.

Notice that $\sym\alpha=0$ for antisymmetric tensors of dimension $>1$ and similarly $\alt\alpha=0$ for symmetric tensors of dimension $>1$.

\subsection{Tensor fields}

\bdefn

    Let $M$ be a smooth manifold.
    Define its {\emphcolor bundle of covariant $k$-tensors} by
    $$ T^kT^*M = \coprod_{p\in M}T^k(T^*_pM) $$
    where $T_p^*M$ is the dual space of $T_pM$.
    Similarly the {\emphcolor bundle of contravariant $k$-tensors} by
    $$ T^kTM = \coprod_{p\in M}T^k(T_pM) , $$
    and in general the {\emphcolor bundle of $(k,\ell)$-tensors} by
    $$ T^{(k,\ell)}TM = \coprod_{p\in M}T^{(k,\ell)}(T_pM) . $$

    $T^*M=T^1T^*M=T^{(0,1}TM$ is also called the {\emphcolor cotangent bundle}.

\edefn


